% ═══════════════════════════════════════════════════════════════════════════
%  EXPERIMENT II — Velocity sweep with Monte Carlo runs (UPDATED real data)
%  Requires: \usepackage{booktabs, siunitx, amsmath, amssymb, graphicx}
% ═══════════════════════════════════════════════════════════════════════════

%  NOTE: The unified file "implementation_and_pareto.tex" contains
%  the canonical version of this section.  This file is the standalone
%  backup with the same corrected numbers.

To assess robustness beyond the single-run baseline, we conduct a
velocity-sweep experiment over three maximum progress speeds
$v_{\theta,\max} \in \{8,\,12,\,15\}\;\si{m/s}$,
spanning conservative to aggressive flight regimes.
For each speed, $N = 5$ Monte Carlo trials are executed with
randomly perturbed initial conditions
($\sigma_p = \SI{0.05}{m}$, $\sigma_q = \SI{0.05}{rad}$,
 seed = 42)
on the same \SI{100}{m} Lissajous path; all other parameters
(Optuna-tuned cost weights, constraints, prediction horizon,
control frequency) remain identical to Experiment~I.
This design isolates the effect of the state representation under
increasing dynamic demand: higher $v_{\theta,\max}$ forces the solver
to plan faster motions, amplifying any shortcomings in the pose
error model.

\begin{table}[t]
  \centering
  \caption{Experiment~II — median RMSE (std) over $N=5$ Monte Carlo
           runs per speed.  $\Delta\%$ is computed on medians.
           Zero failures for both controllers at all speeds.}
  \label{tab:exp2}
  \setlength{\tabcolsep}{3pt}
  \footnotesize
  \begin{tabular}{c cc cc c}
    \toprule
    & \multicolumn{2}{c}{\textbf{DQ-MPCC}}
    & \multicolumn{2}{c}{\textbf{MPCC Baseline}}
    & \\
    \cmidrule(lr){2-3} \cmidrule(lr){4-5}
    $v_{\theta,\max}$
    & {RMSE$_{\text{pos}}$}
    & {RMSE$_{\text{ori}}$}
    & {RMSE$_{\text{pos}}$}
    & {RMSE$_{\text{ori}}$}
    & {$\Delta\%_{\text{pos}}$} \\
    {[m/s]}
    & {[m]}
    & {[rad]}
    & {[m]}
    & {[rad]}
    & \\
    \midrule
    8
    & 0.538\,{\scriptsize($\pm$0.003)}
    & 0.816\,{\scriptsize($\pm$0.001)}
    & 0.673\,{\scriptsize($\pm$0.009)}
    & 0.733\,{\scriptsize($\pm$0.002)}
    & $-20.0\%$ \\
    12
    & 0.623\,{\scriptsize($\pm$0.004)}
    & 0.967\,{\scriptsize($\pm$0.003)}
    & 1.137\,{\scriptsize($\pm$0.146)}
    & 0.969\,{\scriptsize($\pm$0.053)}
    & $-45.2\%$ \\
    15
    & 0.700\,{\scriptsize($\pm$0.006)}
    & 1.050\,{\scriptsize($\pm$0.001)}
    & 1.226\,{\scriptsize($\pm$0.012)}
    & 0.965\,{\scriptsize($\pm$0.005)}
    & $-42.9\%$ \\
    \bottomrule
  \end{tabular}
\end{table}

Three key observations emerge from Table~\ref{tab:exp2}.

\paragraph{DQ-MPCC dominates at every speed.}
At $v_{\theta,\max} = \SI{8}{m/s}$, the improvement in position
RMSE is a moderate~$20\%$ ($\SI{0.538}{m}$ vs.\ $\SI{0.673}{m}$).
At $\SI{12}{m/s}$ and $\SI{15}{m/s}$, the gap widens to
$45\%$ and $43\%$, respectively.
This is consistent with the theoretical advantage of the
$\mathrm{SE}(3)$ logarithmic error: at low speeds the
rotation–translation coupling is mild and both representations
perform comparably, but at high speeds the decoupled
$\mathrm{SO}(3) + \mathbb{R}^3$ error of the baseline fails
to anticipate the interaction between angular and linear
dynamics, leading to larger deviations in the high-curvature
segments of the Lissajous.

\paragraph{DQ-MPCC is far more repeatable.}
The standard deviation of the baseline's RMSE$_{\text{pos}}$
at $v_{\theta,\max} = \SI{12}{m/s}$ is $\pm\SI{0.146}{m}$
(an order of magnitude larger than DQ-MPCC's
$\pm\SI{0.004}{m}$), indicating high sensitivity to the
random initial-condition perturbations.
DQ-MPCC maintains near-deterministic performance across all
five trials, with standard deviations below~\SI{0.006}{m}
at every speed.

\paragraph{Orientation metrics are comparable.}
The baseline achieves slightly lower orientation RMSE at
$v_{\theta,\max} = 8$ and $\SI{15}{m/s}$
($\SI{0.733}{rad}$ vs.\ $\SI{0.816}{rad}$, and
 $\SI{0.965}{rad}$ vs.\ $\SI{1.050}{rad}$).
At $\SI{12}{m/s}$ the two controllers are effectively tied
($\approx\!\SI{0.97}{rad}$).
This marginal orientation penalty is expected:
the $\mathrm{SE}(3)$ logarithm couples rotation and
translation in a single cost term, so the optimiser
allocates part of the rotational budget to reduce
the dominant positional component.

\paragraph{Reconciliation with Experiment~I.}
The single-run comparison (Table~\ref{tab:exp1}) showed
the baseline with lower RMSE and faster $t_{\text{lap}}$
at $v_{\theta,\max} = \SI{15}{m/s}$.
The velocity sweep reveals that this was partly an artefact
of the specific initial condition: the baseline's
\emph{median} RMSE$_{\text{pos}}$ at $\SI{15}{m/s}$
is $\SI{1.226}{m}$ (Table~\ref{tab:exp2}), which is
$78\%$ worse than DQ-MPCC's $\SI{0.700}{m}$ when averaged
over multiple perturbed starts.
The Monte Carlo design eliminates this luck factor and
confirms DQ-MPCC's superiority under statistical rigour.
